% 15-bounds.tex — Chapter 15: Static Bounds Checking
\chapter{Static Bounds Checking}
\label{chap:15-bounds}
\index{bounds checking}
\index{static bounds checking}

\pnum
The bounds checker verifies that every array and slice subscript is within the
valid range $[0, \text{length}-1]$ before allowing the expression to reach the
code generator.  It is implemented in \srcref{src/sema/bounds.h} and called from
the type-checking walk (\cname{sema\_infer\_expr}) when processing
\cname{EXPR\_INDEX} nodes.

% -------------------------------------------------------------------------
\section{Check algorithm}
\label{sec:bounds-algo}
% -------------------------------------------------------------------------

\pnum
For each \cname{EXPR\_INDEX} expression (syntax: \lain{arr[idx]}):

\begin{algorithm}
\textbf{Bounds check for \lain{arr[idx]}}:
\begin{enumerate}
\item If \cname{sema\_in\_unsafe\_block} is true: skip all checks (unsafe block
  explicitly disables bounds verification).
\item If \cname{sema\_walk\_phase} is false: skip (bounds checking only runs
  during walk phase).
\item If \lain{idx} is an \cname{EXPR\_RANGE} (\lain{arr[lo..hi]}): check both
  endpoints independently.
\item \textbf{In-guard check}: if \cname{sema\_is\_in\_guarded(idx, arr)} returns
  true, the index is proven in-bounds by a live \lain{in} guard; skip.
\item \textbf{VRA check}: evaluate \cname{sema\_eval\_range(idx, sema\_ranges)}.
  If the resulting range $[min, max]$ satisfies $0 \leq min$ and $max < arr.len$
  (for a statically known length), the access is proven safe; skip.
\item \textbf{Runtime check}: if neither static proof applies and the array length
  is known, emit a runtime bounds check in the generated C output.
\item If the VRA proves the access is \emph{definitely out of bounds} (i.e.\
  $min \geq arr.len$ or $max < 0$), emit a compile-time error \diagcode{E045}
  (or an appropriate bounds error code).
\end{enumerate}
\end{algorithm}

% -------------------------------------------------------------------------
\section{Safe fragment guarantee}
\label{sec:bounds-safe}
% -------------------------------------------------------------------------

\pnum
In the safe fragment (no \lain{unsafe} blocks), the compiler guarantees that
all array accesses are either statically proven in-bounds or guarded by a
runtime check in the generated C code.  This guarantee is contingent on:

\begin{itemize}
\item The VRA not being circumvented by pointer arithmetic.
\item The in-guard table being correctly maintained (no double-pop errors).
\item The \cname{sema\_walk\_phase} flag being set correctly during walk.
\end{itemize}

\begin{invariant}
No \cname{EXPR\_INDEX} node that reaches \cname{emit\_expr()} has been
subjected to bounds checking in \lain{unsafe} mode without explicit opt-in.
Every subscript either has a static proof or a generated runtime check.
\end{invariant}

% -------------------------------------------------------------------------
\section{Interaction with VRA}
\label{sec:bounds-vra}
% -------------------------------------------------------------------------

\pnum
The bounds checker cooperates with VRA in two directions:
\begin{enumerate}
\item \textbf{VRA proves bounds}: if VRA has refined the range of the index
  to be within $[0, len-1]$, no runtime check is emitted.
\item \textbf{Constraint seeding}: \lain{in} conditions (\lain{idx in arr})
  seed the VRA with the constraint $0 \leq idx < arr.len$, which is then
  available as a proven fact for all code inside the guard.
\end{enumerate}

% -------------------------------------------------------------------------
\section{Interaction with unsafe}
\label{sec:bounds-unsafe}
% -------------------------------------------------------------------------

\pnum
Inside \lain{unsafe \{ ... \}} blocks, \cname{sema\_in\_unsafe\_block} is
\texttt{true} and bounds checking is skipped entirely.  The emitter generates
raw unchecked array subscripts.  The programmer is responsible for any
out-of-bounds accesses that may result.  This is the intended escape hatch for
performance-critical code that has been verified by other means.
